\scanpage{32}
\uline{NB}

\[
\begin{aligned}
\overset{\substack{C_Q\\\rotatebox[origin=c]{90}{$\scriptstyle\simeq$}}}{\underline\omega}&\subset\mathbb F_2^{\underline\omega}=\mathbb F_2^C=V_S\subset\mathcal D_Q&
\underline\omega&=V_S\setminus\mathfrak z_Q&\longleftrightarrow&\ \mathrm{sg}(g)\\
D_Q&\simeq\underline\omega\wedge\underline\mu\subset\mathbb F_2^D=V_A\subset\mathcal D_Q&
\underline\omega\wedge\mu&\simeq V_A\setminus\mathfrak z_Q&\longleftrightarrow&\ \mathrm{sg}(g)\det(g)\\
\underline\omega_Q&\simeq\underline\mu\subset\mathcal D_Q^+\subset\mathcal D_Q&
\underline\mu&\simeq\mathcal D_Q^+\setminus\mathfrak z_Q&\longleftrightarrow&\ \det(g)
\end{aligned}
\]

\[
\begin{aligned}
\text{où}\quad \underline\mu&\unsure{=|\simeq}(\overset{2}{\wedge} M)^*=\mathrm{Or}(M)\\
\mathfrak z_Q&=\{1,\underline a_Q\}=\mathrm{Cent}\,\mathcal D_Q
\end{aligned}
\]

\margin{$\mathrm{Gl}_{\Lambda}(M)'\simeq\mathfrak S_I\times\mathbb F_2\simeq\mathrm{Aut}(\Delta)$ où $\Delta$ est l'octaèdre non orienté associé à $M$ de sommets $\Delta=M^{\unsure{+|*|\#}}/\{\pm1\}$ \ldots}

\diagram{\scriptsize\begin{tikzcd}[ampersand replacement=\&,column sep=0.1em,row sep=1.8em,cramped,execute at end picture={
\draw ([xshift=-4pt]fw.north west) to[out=65,in=230] ($(fw.north west)!0.5!(top.north west)+(-10pt,6pt)$) to[out=60,in=200] ([xshift=-2pt,yshift=5pt]top.north west);
\node[above left] at ($(fw.north west)!0.5!(top.north west)+(-10pt,6pt)$) {$\mathfrak S_{\unsure{I|J}}$};
\draw (vec.north west) to[out=75,in=235] ($(vec.north west)!0.5!(top.west)+(-3pt,0pt)$) to[out=65,in=205] (top.west);
\node[above left,align=center] at ($(vec.north west)!0.5!(top.west)$) {$\mathfrak S_J$\\$\simeq\mathrm{Gl}_{\mathbb F_2}(M_0)$};
\draw ([xshift=4pt]mo.south east) to[out=10,in=230] ($(mo.east)!0.5!(alt.east)+(5pt,-2pt)$) to[out=90,in=245] (alt.south east);
\node[right] at ($(mo.east)!0.5!(alt.east)+(5pt,-2pt)$) {$\mathfrak S_J^+$};
\draw (sl.south east) to[out=-85,in=95] ($(sl.east)!0.5!(alt.east)+(8pt,0pt)$) to[out=-95,in=85] (alt.east);
\node[right,align=left] (quot) at ($(sl.east)!0.5!(alt.east)+(8pt,0pt)$) {$\mathcal D_Q^+\simeq(\mathrm{Sl}_{\Lambda}(M))_{\mathrm{ab}}$\\$\simeq\overset{2}{\wedge}_{\deleted{\unsure{\mathbb Z/4\mathbb Z|\Lambda}}} M$};
\draw (quot.east |- top.north) to[out=-88,in=92] ([xshift=3pt]quot.north east) to[out=-90,in=84] (quot.east |- alt.south);
\node[right] at (quot.east |- sl.south) {$\mathcal D_Q$};
}]
\&\&\&|[alias=top]|{\mathrm{Gl}_{\Lambda}(M)[\simeq\mathfrak S_I\cdot\mathbb F_2^{\underline\omega}]}\&\\
\&\&\mathfrak S_I^+\times\mathbb F_2^{\underline\omega}\arrow[ur,"2"]\&{[\mathfrak S_I\times\mathfrak z]}\arrow[u,dotted]\&|[alias=sl]|{\mathrm{Sl}_{\Lambda}(M)}\arrow[ul,"2"']\\
\&|[alias=vec]|{M_0\times\mathbb F_2^{\underline\omega}\simeq\mathfrak{gl}(M_0)}\arrow[ur,"3"]\&\&\mathfrak S_I^+\times\mathfrak z\arrow[ul,"2"]\arrow[u,dotted]\arrow[ur,"2"]\&\\
|[alias=fw]|{\mathbb F_2^{\underline\omega}}\arrow[ur,"4"]\&\&M_0\times\mathfrak z\simeq\mathfrak{sl}(M_0)\arrow[ul,"2"]\arrow[ur,"3"]\&\&|[alias=alt]|{\mathfrak S_I^+}\arrow[ul,"2"]\\
{(28)}\&\mathfrak z\arrow[ul,"2"]\arrow[ur,"4"]\&\&|[alias=mo]|{M_0}\arrow[ul,"2"]\arrow[ur,"3"]\&\\
\&\&1\arrow[ul,"2"]\arrow[ur,"4"]\&\&
\end{tikzcd}}

\uline{NB} Dans ce diagramme, toutes les flèches sont injectives, et les cinq carrés qui y apparaissent, donc aussi les composés, sont cartésiens. On a marqué sur chaque flèche l'indice \deleted{\unsure{pour}} \unsure{de|pour} l'inclusion de sous-groupes qu'elle représente.

Les iso.

\[
\begin{aligned}
\mathrm{Gl}_{\Lambda}(M)&\simeq\mathfrak S_I\cdot\mathbb F_2^{\underline\omega}\quad(1/2\text{ direct})\quad\text{et}\\
\mathrm{Ker}(\mathrm{sg}\cdot\det)&\simeq\mathfrak S_I\times\deleted{\mathbb F_2}\mathfrak z\quad(\text{prod direct}
\end{aligned}
\]

sont valables quand $M$ est \textquotedblleft $1/2$ déployé\textquotedblright\ i.e. le torseur $P_M$ est trivialisé, alors on a $=$

\[
\mathrm{Gl}_{\Lambda}(M)\simeq
\underbrace{\mathrm{Gl}_{\mathbb F_2}(M_0)}_{\mathfrak S_J}\cdot(M_0\times\mathbb F_2^{\underline\omega})
\simeq\underbrace{(\mathrm{Gl}_{\mathbb F_2}(M_0)\cdot M_0)}_{\mathfrak S_I}\cdot\mathbb F_2^{\underline\omega}
\]

mais ces isom. dépendent du choix d'une trivialisat. de $P_M$ \ldots
% archive sheet number: 31
